% ============================================================
% 8.定积分.tex —— 第5章 定积分（同济《高等数学》第七版·上册）
% 第一节 定积分的概念与性质（问题举例 / 定义 / 近似计算 / 性质）
% 第二节 微积分基本公式（位置与速度 / 积分上限函数 / 牛顿—莱布尼茨公式）
% 第三节 定积分的换元法和分部积分法
% 第四节 反常积分（无穷限 / 无界函数）
% 第五节 反常积分的审敛法 Γ函数
% 只写知识点，不抄例题与习题。
% ============================================================

\section{定积分的概念与性质}

\subsection{定积分问题举例}

\noindent\textbf{曲边梯形的面积}\quad 设$f(x)\in C[a,b]$，$f(x)\ge0$。由$y=f(x)$、$x$轴及$x=a$、$x=b$围成的图形称为\emph{曲边梯形}。求面积$A$的四步：
\begin{enumerate}
  \item \textbf{分割}：以分点$a=x_0<x_1<\cdots<x_n=b$把$[a,b]$分成$n$个小区间，第$i$个窄曲边梯形面积记为$\Delta A_i$；
  \item \textbf{近似}：任取$\xi_i\in[x_{i-1},x_i]$，以$f(\xi_i)\Delta x_i$近似代替$\Delta A_i$（以直代曲、以常代变）；
  \item \textbf{求和}：$A\approx\sum\limits_{i=1}^{n}f(\xi_i)\Delta x_i$；
  \item \textbf{取极限}：令$\lambda=\max\{\Delta x_i\}\to0$，得$A=\lim\limits_{\lambda\to0}\sum\limits_{i=1}^{n}f(\xi_i)\Delta x_i$。
\end{enumerate}

\noindent\textbf{变速直线运动的路程}\quad 设$v=v(t)\ge0$连续，则$s=\lim\limits_{\lambda\to0}\sum\limits_{i=1}^{n}v(\tau_i)\Delta t_i$。两个问题意义不同，却归结为同一数学结构，由此抽象出定积分。

\subsection{定积分的定义}

\noindent\textbf{定义5.1}\quad 设$f(x)$在$[a,b]$上有界，用分点
\[
a=x_0<x_1<x_2<\cdots<x_{i-1}<x_i<\cdots<x_n=b
\]
把$[a,b]$分成$n$个小区间$[x_{i-1},x_i]$，记$\Delta x_i=x_i-x_{i-1}$，$\lambda=\max\limits_{1\le i\le n}\{\Delta x_i\}$。在每个小区间上任取$\xi_i\in[x_{i-1},x_i]$，作\emph{积分和}
$S_n=\sum\limits_{i=1}^{n}f(\xi_i)\Delta x_i$。如果当$\lambda\to0$时$S_n$的极限存在，且\textbf{与分法无关、与$\xi_i$的取法无关}，则称$f(x)$在$[a,b]$上\emph{可积}，并称此极限为$f(x)$在$[a,b]$上的\emph{定积分}：
\[
\int_a^b f(x)\,\mathrm{d}x=\lim_{\lambda\to0}\sum_{i=1}^{n}f(\xi_i)\Delta x_i .
\]
其中$a$为积分下限，$b$为积分上限，$[a,b]$为积分区间。

\keyformula{\int_a^b f(x)\,\mathrm{d}x=\lim_{\lambda\to0}\sum_{i=1}^{n}f(\xi_i)\Delta x_i}

\noindent\textbf{规定}\quad $\displaystyle\int_a^af(x)\mathrm{d}x=0$，$\displaystyle\int_a^bf(x)\mathrm{d}x=-\int_b^af(x)\mathrm{d}x\ (a>b)$。

\noindent\textbf{几何意义}\quad $f\ge0$时为曲边梯形面积；$f\le0$时为面积的相反数；$f$变号时为$x$轴上、下方面积的\textbf{代数和}。

\noindent\textbf{可积性}\quad (1) \textbf{必要}：可积必有界（有界不充分）；(2) \textbf{充分}：$f\in C[a,b]$；(3) \textbf{充分}：$f$有界且只有\emph{有限个间断点}。

\corepoint{定积分是一个\textbf{数}，只与区间和被积函数有关，与积分变量字母无关：$\int_a^bf(x)\mathrm{d}x=\int_a^bf(t)\mathrm{d}t$。$\mathrm{d}x$必须写，它标明积分变量；结果\textbf{不带$C$}。}

\subsection{定积分的近似计算}

把$[a,b]$分成$n$等份，$h=\dfrac{b-a}{n}$，$x_i=a+ih$，则

\noindent\textbf{矩形法}\quad
$\displaystyle\int_a^bf(x)\mathrm{d}x\approx\frac{b-a}{n}\sum_{i=1}^{n}f(x_{i-1})$（左端点）或$\displaystyle\approx\frac{b-a}{n}\sum_{i=1}^{n}f(x_i)$（右端点）。

\noindent\textbf{梯形法}\quad 以折线代曲线：
\[
\int_a^bf(x)\,\mathrm{d}x\approx\frac{b-a}{n}\left[\frac{f(a)+f(b)}{2}+\sum_{i=1}^{n-1}f(x_i)\right].
\]

\noindent\textbf{抛物线法（辛普森公式）}\quad 要求$n$为\textbf{偶数}，过每三个相邻分点作抛物线：
\[
\int_a^bf(x)\,\mathrm{d}x\approx\frac{b-a}{3n}\left[f(a)+f(b)+4\!\!\sum_{i=1}^{n/2}f(x_{2i-1})+2\!\!\sum_{i=1}^{n/2-1}f(x_{2i})\right].
\]

\corepoint{三种近似精度递增：矩形法最差，梯形法次之；抛物线法对次数$\le3$的多项式\textbf{精确成立}，同分点下最准。$n$必须为偶数才能用抛物线法，这是最易漏的前提。}

\subsection{定积分的性质}

以下设所涉积分均存在。

\noindent\textbf{性质1（线性）}\quad $\displaystyle\int_a^b[kf+lg]\mathrm{d}x=k\int_a^bf\mathrm{d}x+l\int_a^bg\mathrm{d}x$。

\noindent\textbf{性质2（区间可加）}\quad 对任意$a,b,c$（不要求$a<c<b$），$\displaystyle\int_a^bf\mathrm{d}x=\int_a^cf\mathrm{d}x+\int_c^bf\mathrm{d}x$。

\noindent\textbf{性质3（保号性）}\quad $f\ge0$且$a<b$ $\Rightarrow$ $\displaystyle\int_a^bf\mathrm{d}x\ge0$。

\noindent\textbf{性质4（保序性）}\quad $f\le g$ $\Rightarrow$ $\displaystyle\int_a^bf\mathrm{d}x\le\int_a^bg\mathrm{d}x$。

\noindent\textbf{性质5（绝对值不等式）}\quad $\left|\displaystyle\int_a^bf\mathrm{d}x\right|\le\displaystyle\int_a^b|f|\mathrm{d}x$。

\noindent\textbf{性质6（估值定理）}\quad 若$m\le f(x)\le M$（$a<b$），则
\[
m(b-a)\le\int_a^bf(x)\,\mathrm{d}x\le M(b-a).
\]

\noindent\textbf{性质7（积分中值定理）}\quad 若$f\in C[a,b]$，则至少存在$\xi\in[a,b]$使
\[
\int_a^bf(x)\,\mathrm{d}x=f(\xi)(b-a),
\]
而$f(\xi)=\dfrac{1}{b-a}\displaystyle\int_a^bf\,\mathrm{d}x$称为$f$在$[a,b]$上的\emph{平均值}。

\keyformula{m(b-a)\le\int_a^bf\mathrm{d}x\le M(b-a);\qquad \int_a^bf\mathrm{d}x=f(\xi)(b-a)}

\corepoint{保号性要求\emph{下限小于上限}。$f$连续且$\int_a^bf\mathrm{d}x=0$只说明$f$在$[a,b]$上\textbf{变号或为零}（至少一个零点），加$f\ge0$才能推出$f\equiv0$。中值定理的$\xi$只保证存在，不能当具体值用。}

\section{微积分基本公式}

\subsection{变速直线运动中位置函数与速度函数之间的联系}

设位置$s=s(t)$，速度$v=v(t)=s'(t)$，则区间$[T_1,T_2]$上的路程既等于$\int_{T_1}^{T_2}v\,\mathrm{d}t$，又等于$s(T_2)-s(T_1)$，故
\[
\int_{T_1}^{T_2}v(t)\,\mathrm{d}t=s(T_2)-s(T_1),
\]
即定积分可用被积函数的原函数在端点处的值之差来计算。此事实对一般连续函数也成立。

\subsection{积分上限的函数及其导数}

设$f\in C[a,b]$，定义$\Phi(x)=\displaystyle\int_a^xf(t)\,\mathrm{d}t$，称为\emph{积分上限函数}（变上限积分）。

\noindent\textbf{定理5.1（原函数存在定理）}\quad 若$f\in C[a,b]$，则$\Phi(x)$在$[a,b]$上可导，且
\[
\Phi'(x)=\frac{\mathrm{d}}{\mathrm{d}x}\int_a^xf(t)\,\mathrm{d}t=f(x).
\]
即\textbf{连续函数的变上限积分是它的一个原函数}，从而连续函数必有原函数。

\noindent\textbf{定理5.2（复合形式）}\quad 设$f$连续，$\varphi,\psi$可导，则
\[
\frac{\mathrm{d}}{\mathrm{d}x}\int_{\varphi(x)}^{\psi(x)}f(t)\,\mathrm{d}t=f[\psi(x)]\psi'(x)-f[\varphi(x)]\varphi'(x).
\]

\keyformula{\frac{\mathrm{d}}{\mathrm{d}x}\int_a^xf(t)\mathrm{d}t=f(x);\qquad \frac{\mathrm{d}}{\mathrm{d}x}\int_{\varphi(x)}^{\psi(x)}f(t)\mathrm{d}t=f(\psi)\psi'-f(\varphi)\varphi'}

\corepoint{变上限求导时被积函数\textbf{不能含$x$}；形如$\int_a^xf(x-t)\mathrm{d}t$必须先换元（令$u=x-t$）把$x$移出被积函数。上下限都是函数时，别忘"上限导数$-$下限导数"的减号。}

\subsection{牛顿—莱布尼茨公式}

\noindent\textbf{定理5.3（微积分基本定理）}\quad 若$F(x)$是连续函数$f(x)$在$[a,b]$上的一个原函数，则
\[
\int_a^bf(x)\,\mathrm{d}x=F(b)-F(a)=\bigl[F(x)\bigr]_a^b ,
\]
称为\emph{牛顿—莱布尼茨公式}（微积分基本公式）。

\noindent\textbf{意义}\quad 把求定积分（积分和的极限）归结为求原函数，使定积分可算；同时揭示微分与积分互为逆运算。

\noindent\textbf{步骤}\quad (1) 检验$f$在$[a,b]$上连续；(2) 求$f$的\emph{一个}原函数$F$（不带$C$）；(3) 代入上下限作差$F(b)-F(a)$。

\corepoint{公式条件是$f$在$[a,b]$上\textbf{连续}。若区间内有瑕点（如$\int_{-1}^1\frac{\mathrm{d}x}{x^2}$），\textbf{不能}直接用N-L公式，须按反常积分处理；硬套会得出$-2$之类的荒谬结果。}

\section{定积分的换元法和分部积分法}

\subsection{定积分的换元法}

\noindent\textbf{定理5.4}\quad 设$f\in C[a,b]$，$x=\varphi(t)$满足：(1) $\varphi(\alpha)=a$，$\varphi(\beta)=b$，且$t$在$\alpha,\beta$之间时$\varphi(t)\in[a,b]$；(2) $\varphi'(t)$在$[\alpha,\beta]$（或$[\beta,\alpha]$）上连续。则
\[
\int_a^bf(x)\,\mathrm{d}x=\int_{\alpha}^{\beta}f[\varphi(t)]\varphi'(t)\,\mathrm{d}t .
\]

\noindent\textbf{与不定积分的区别}\quad 定积分换元\textbf{必须同时换积分限}，且\textbf{不需回代}：求出关于$t$的原函数后直接用$t$的上下限作差。

\noindent\textbf{对称区间的奇偶性}\quad 设$f$在$[-a,a]$上连续，则
\[
\int_{-a}^{a}f(x)\,\mathrm{d}x=
\begin{cases}
0, & f(-x)=-f(x)\ (\text{奇}),\\[4pt]
2\displaystyle\int_0^{a}f(x)\,\mathrm{d}x, & f(-x)=f(x)\ (\text{偶}).
\end{cases}
\]

\noindent\textbf{常用结论}（$f$连续）
\begin{align*}
&\int_0^{\pi}xf(\sin x)\,\mathrm{d}x=\frac{\pi}{2}\int_0^{\pi}f(\sin x)\,\mathrm{d}x,\qquad
\int_0^{\pi/2}f(\sin x)\,\mathrm{d}x=\int_0^{\pi/2}f(\cos x)\,\mathrm{d}x,\\
&\int_0^{\pi/2}\sin^nx\,\mathrm{d}x=\int_0^{\pi/2}\cos^nx\,\mathrm{d}x=
\begin{cases}
\dfrac{n-1}{n}\cdot\dfrac{n-3}{n-2}\cdots\dfrac{1}{2}\cdot\dfrac{\pi}{2}, & n\ \text{偶},\\[8pt]
\dfrac{n-1}{n}\cdot\dfrac{n-3}{n-2}\cdots\dfrac{2}{3}, & n\ \text{奇}.
\end{cases}
\end{align*}

\corepoint{三条铁律：换元必换限；限的对应由$x=\varphi(t)$直接算出；$\varphi(t)$须单调（教材条件已保证），否则一个$x$对应多个$t$会算重。用奇偶性时区间必须\textbf{关于原点对称}。}

\subsection{定积分的分部积分法}

\noindent\textbf{定理5.5}\quad 设$u(x),v(x)$在$[a,b]$上有连续导数，则
\[
\int_a^bu\,\mathrm{d}v=\bigl[uv\bigr]_a^b-\int_a^bv\,\mathrm{d}u ,
\quad\text{即}\quad
\int_a^bu(x)v'(x)\,\mathrm{d}x=\bigl[u(x)v(x)\bigr]_a^b-\int_a^bu'(x)v(x)\,\mathrm{d}x .
\]
$u$的选取顺序仍为"反对幂指三"。记$I_n=\displaystyle\int_0^{\pi/2}\sin^nx\,\mathrm{d}x$，则$I_n=\dfrac{n-1}{n}I_{n-2}$（$n\ge2$），$I_0=\dfrac{\pi}{2}$，$I_1=1$。

\corepoint{$\bigl[uv\bigr]_a^b$这一项\textbf{必须代上下限}，忘了等于只做了不定积分。被积函数含绝对值或分段定义时，先在分界点\textbf{拆区间}再分部。}

\section{反常积分}

定积分$\int_a^bf\mathrm{d}x$有两个前提：区间\emph{有限}、被积函数\emph{有界}。突破其一即为\emph{反常积分}（广义积分）。

\subsection{无穷限的反常积分}

\noindent\textbf{定义5.2}\quad 设$f$在$[a,+\infty)$上连续，$b>a$。若$\lim\limits_{b\to+\infty}\displaystyle\int_a^bf(x)\mathrm{d}x$存在，则称此极限为$f$在$[a,+\infty)$上的反常积分，记
\[
\int_a^{+\infty}f(x)\,\mathrm{d}x=\lim_{b\to+\infty}\int_a^bf(x)\,\mathrm{d}x ,
\]
并称它\emph{收敛}；极限不存在则称\emph{发散}。类似地定义$\displaystyle\int_{-\infty}^bf\mathrm{d}x$，且
\[
\int_{-\infty}^{+\infty}f(x)\,\mathrm{d}x=\int_{-\infty}^{c}f(x)\,\mathrm{d}x+\int_c^{+\infty}f(x)\,\mathrm{d}x
\]
当右端\textbf{两个积分都收敛}时才收敛（$c$任取）。

\noindent\textbf{记法}\quad 令$F(+\infty)=\lim\limits_{x\to+\infty}F(x)$，$F(-\infty)=\lim\limits_{x\to-\infty}F(x)$，则$\displaystyle\int_a^{+\infty}f\mathrm{d}x=F(x)\big|_a^{+\infty}=F(+\infty)-F(a)$；收敛性取决于$F(\pm\infty)$是否为有限值。

\noindent\textbf{$p$积分}\quad $\displaystyle\int_a^{+\infty}\frac{\mathrm{d}x}{x^p}\ (a>0)$：$p>1$收敛，$p\le1$发散。

\corepoint{$\int_{-\infty}^{+\infty}f$的收敛是"两段各自收敛"，\textbf{不是}主值$\lim\limits_{R\to+\infty}\int_{-R}^{R}f$。例如$\int_{-\infty}^{+\infty}x\mathrm{d}x$发散（主值为0）。必须\textbf{拆两段}分别算。}

\subsection{无界函数的反常积分}

\noindent\textbf{定义5.3}\quad 设$f$在$(a,b]$上连续，在$a$的右邻域内无界（$a$称为\emph{瑕点}）。若$\lim\limits_{\varepsilon\to0^+}\displaystyle\int_{a+\varepsilon}^bf(x)\mathrm{d}x$存在，则称此极限为$f$在$(a,b]$上的反常积分（\emph{瑕积分}），仍记$\displaystyle\int_a^bf(x)\mathrm{d}x$，即
\[
\int_a^bf(x)\,\mathrm{d}x=\lim_{\varepsilon\to0^+}\int_{a+\varepsilon}^{b}f(x)\,\mathrm{d}x ,
\]
并称它收敛，否则发散。瑕点为$b$时取$\lim\limits_{\varepsilon\to0^+}\int_a^{b-\varepsilon}f\mathrm{d}x$；瑕点$c\in(a,b)$时须\textbf{两端都收敛}。

\noindent\textbf{瑕积分$p$判别}\quad $\displaystyle\int_a^b\frac{\mathrm{d}x}{(x-a)^p}$：$p<1$收敛，$p\ge1$发散（$b>a$）。

\corepoint{先\textbf{找瑕点}（使被积函数趋于无穷的点）。常漏判：$\int_0^1\frac{\mathrm{d}x}{\sqrt{x(1-x)}}$（两端皆瑕点）、$\int_0^{\pi}\frac{\mathrm{d}x}{\sin x}$（$x=0,\pi$）、$\int_0^1\frac{\ln x}{x}\mathrm{d}x$。有瑕点就不能直接套N-L公式。}

\section{反常积分的审敛法 \texorpdfstring{$\Gamma$}{Gamma}函数}

\subsection{无穷限反常积分的审敛法}

总设$f$在$[a,+\infty)$上连续且$f(x)\ge0$。

\noindent\textbf{定理5.6（比较判别法）}\quad 设$0\le f\le g\ (x\ge a)$：(1) 若$\int_a^{+\infty}g$收敛，则$\int_a^{+\infty}f$收敛；(2) 若$\int_a^{+\infty}f$发散，则$\int_a^{+\infty}g$发散。

\noindent\textbf{定理5.7（极限判别法）}\quad (1) 若存在$p>1$使$\lim\limits_{x\to+\infty}x^pf(x)=c\ (0\le c<+\infty)$，则$\int_a^{+\infty}f$收敛；(2) 若存在$p\le1$使$\lim\limits_{x\to+\infty}x^pf(x)=c\ (0<c\le+\infty)$，则$\int_a^{+\infty}f$发散。

\noindent\textbf{定理5.8（绝对收敛）}\quad 若$\int_a^{+\infty}|f|\mathrm{d}x$收敛，则$\int_a^{+\infty}f\mathrm{d}x$也收敛，称\emph{绝对收敛}；若$\int f$收敛而$\int|f|$发散，称\emph{条件收敛}（如$\int_1^{+\infty}\frac{\sin x}{x}\mathrm{d}x$）。

\subsection{无界函数的反常积分的审敛法}

设$x=a$为瑕点，$f(x)\ge0$，$f$在$(a,b]$上连续。

\noindent\textbf{比较判别法}\quad 设$0\le f\le g$：(1) 若$\int_a^bg$收敛，则$\int_a^bf$收敛；(2) 若$\int_a^bf$发散，则$\int_a^bg$发散。

\noindent\textbf{极限判别法}\quad (1) 若存在$p<1$使$\lim\limits_{x\to a^+}(x-a)^pf(x)=c\ (0\le c<+\infty)$，则收敛；(2) 若存在$p\ge1$使$\lim\limits_{x\to a^+}(x-a)^pf(x)=c\ (0<c\le+\infty)$，则发散。

\keyformula{\int_a^{+\infty}\frac{\mathrm{d}x}{x^p}:p>1\ \text{收敛};\qquad \int_a^b\frac{\mathrm{d}x}{(x-a)^p}:p<1\ \text{收敛}}

\corepoint{无穷限与瑕积分的$p$条件\textbf{相反}（一个要$p>1$，一个要$p<1$）。取$p=1$最方便：若$\lim(x-a)f(x)=c>0$即发散。多个瑕点或"既有瑕点又无穷远"时，必须逐段拆开判，全部收敛才算收敛。}

\subsection{\texorpdfstring{$\Gamma$}{Gamma}函数}

\noindent\textbf{定义5.4}\quad 含参变$s$的反常积分
\[
\Gamma(s)=\int_0^{+\infty}x^{s-1}\mathrm{e}^{-x}\,\mathrm{d}x\qquad(s>0)
\]
称为\emph{$\Gamma$函数}（伽马函数）。它在$s>0$时收敛，且在该区间内连续、可导。

\noindent\textbf{性质}
\begin{align*}
&\Gamma(s+1)=s\Gamma(s)\ (s>0);\qquad \Gamma(n+1)=n!\ (n=1,2,\cdots);\\
&\Gamma(1)=1,\qquad \Gamma\!\left(\tfrac12\right)=\sqrt{\pi}.
\end{align*}

\noindent\textbf{常用变形}\quad 令$x=t^2$得$\Gamma(s)=2\displaystyle\int_0^{+\infty}t^{2s-1}\mathrm{e}^{-t^2}\mathrm{d}t$，由此$\displaystyle\int_0^{+\infty}\mathrm{e}^{-x^2}\mathrm{d}x=\frac{\sqrt{\pi}}{2}$，$\displaystyle\int_0^{+\infty}x^2\mathrm{e}^{-x^2}\mathrm{d}x=\frac{\sqrt{\pi}}{4}$。

\corepoint{$\Gamma$的定义域是$s>0$，在$s=0$处发散（$\lim\limits_{s\to0^+}\Gamma(s)=+\infty$）。用递推公式把$s$往$0$推时，每一步都必须保持自变量$>0$。$\Gamma(\frac12)=\sqrt\pi$必背，多数概率积分靠它。}

\begin{center}
\small
\begin{tabular}{@{}lll@{}}
\toprule
\textbf{对象} & \textbf{收敛条件} & \textbf{发散条件} \\
\midrule
$\int_a^{+\infty}\frac{\mathrm{d}x}{x^p}$ & $p>1$ & $p\le1$ \\
$\int_a^b\frac{\mathrm{d}x}{(x-a)^p}$ & $p<1$ & $p\ge1$ \\
\bottomrule
\end{tabular}
\end{center}
